% Psychoacoustic partial-loudness / masked-threshold estimate table.
% Requires packages already in main.tex's preamble: booktabs, siunitx.
% Data source: analysis/compute_partial_loudness.py, run against the same
% 9 exports in "UnMasking Exports/" as parameter_sweep_table.tex. Row order
% matches tab:param-sweep (ranked by mean margin gain) for direct
% cross-reference between the two tables.
%
% METHODOLOGY (condensed - see compute_partial_loudness.py's module
% docstring for the full version):
%  - ISO 532-1 (Zwicker) loudness model via MOSQITO, STATIONARY per-segment
%    method (not the full time-varying model - a computational-cost
%    tradeoff), nperseg=8192/noverlap=4096 @ 48kHz (~85ms hop).
%  - Calibration: full-scale digital peak assumed = 100 dB SPL. Absolute
%    percentages are calibration-dependent; the relative comparison across
%    conditions (same calibration throughout) should be much more robust.
%  - "Masked" = Dialogue's specific loudness in a critical band/frame is
%    <= the reconstructed masker's specific loudness in that same
%    band/frame - the masker's own excitation used as an approximate
%    masking-threshold proxy. This is NOT the full Moore & Glasberg
%    partial-loudness/masking-index model.
%  - Masker reconstructed via exact time-domain subtraction (mix - dry
%    Dialogue), valid because Dialogue (the key/priority channel) always
%    takes the dry, unducked path regardless of duck mode in this engine.
%  - Averaged over frames where Dialogue's own overall loudness exceeds
%    10% of its peak.

\begin{table}[h]
\centering
\footnotesize
\caption{Psychoacoustic partial-loudness estimate: percentage of Dialogue's
own specific loudness (ISO 532-1 Zwicker loudness, stationary per-segment)
falling in critical bands where the reconstructed masker meets or exceeds
it, averaged over active frames. Reduction is in percentage points,
unprocessed minus Advanced. Calibrated so that full-scale digital peak
corresponds to \SI{100}{\decibel} SPL. Row order matches
Table~\ref{tab:param-sweep} for direct comparison.}
\label{tab:partial-loudness}
\begin{tabular}{@{}c S[table-format=2.1] S[table-format=2.1] S[table-format=2.1] S[table-format=+2.1]@{}}
\toprule
\textbf{\#} & {\textbf{Unprocessed} (\%)} & {\textbf{Basic} (\%)} & {\textbf{Advanced} (\%)} & {\textbf{Reduction} (pp)} \\
\midrule
1 & 53.5 & 36.3 & 42.3 & +11.2 \\
2 & 53.5 & 34.1 & 41.3 & +12.2 \\
3 & 53.5 & 39.1 & 44.3 & +9.2 \\
4 & 53.5 & 46.2 & 48.7 & +4.8 \\
5 & 53.5 & 45.2 & 48.1 & +5.4 \\
6 & 53.5 & 45.6 & 48.4 & +5.1 \\
7 & 53.5 & 50.1 & 51.3 & +2.1 \\
8 & 53.5 & 50.3 & 51.5 & +2.0 \\
9 & 53.5 & 52.2 & 52.5 & +0.9 \\
\bottomrule
\end{tabular}
\par\vspace{4pt}
\raggedright\footnotesize Masking-threshold proxy: Dialogue's specific
loudness is treated as masked in a critical band/frame when it does not
exceed the reconstructed masker's specific loudness there. This is an
approximation, not the full Moore \& Glasberg partial-loudness model --
see \texttt{analysis/compute\_partial\_loudness.py} for complete
methodology and caveats.
\end{table}
